\chapter{Modell}\label{sec:modell}

Das in dieser Arbeit verwendete Modell basiert auf dem Nächste-Nachbar-TB-Modell aus Gleichung~\eqref{eq:TB} und der BCS-Wechselwirkung aus Gleichung~\eqref{eq:BCS}. 
Diese werden so modifiziert, dass die Untergitterstruktur des Bienenwabengitters berücksichtigt wird.
Zusätzlich wird noch ein Term zur Berücksichtigung eines chemischer Potentials $\mu$ hinzugefügt, womit der modifizierte Hamiltonian dann lautet:
\begin{align}
    \oH = \oHNN + \oH_{\mu} + \oHWW 
\end{align}
mit 
\begin{align}
    \oHNN &= - t \sum_{\expval{\vec{R}, \vec{R'}}, \sigma} \sum_{\oC, \oC'}\ocd{\vec{R}, \sigma} \op{C'}^{\pdagg}_{\vec{R'}, \sigma}\, , \\
    \oH_{\mu} &= - \mu \sum_{\vec{R}, \sigma} \sum_C \op{C}^{\dagg}_{\vec{R}, \sigma} \op{C}^{\pdagg}_{\vec{R}, \sigma} \, ,\\
    \oHWW &=  - \frac{U}{N} \sump{\vk, \vk'} \sum_\oC  \ocdu \ocdd \ocndd \ocndu \, .
\end{align}

Dabei werden die allgemeinen fermionischen Operatoren $\ofe{}/\ofed{}$  durch die Untergitter-spezifischen Operatoren $\oa{}/\oad{}$ bzw.\ $\ob{}/\obd{}$ ersetzt, die nur Elektronen auf dem jeweiligen Untergitter $A$ oder $B$ erzeugen, bzw.~vernichten. 
Die Summe über $\oC$ bedeutet, dass jeweils über $\oC = \oA$ und $\oC = \oB$ summiert wird. 
% Es fällt auf, dass 

Das chemische Potential $\mu$ kann hier aufgrund der vorhandenen Untergittersymmetrie des Bienenwabengitters und der Spinunabhängigkeit des Systems durch Abwesenheit von externen Magnetfeldern und Vernachlässigung anderer spinabhängiger Effekte als unabhängig von Untergitter und Spin betrachtet werden. 
% Es wird eingeführt um die Wirkung der Verschiebung der Fermienergie durch beispielsweise Dotierung auf die Supraleitung untersuchen zu können.
Es soll die Effekte einer Verschiebung der zuvor auf null gesetzten Fermienergie, durch beispielsweise Dotierung beschreiben.

Da im Bienenwabengitter der nächste Nachbar von jedem $A$-Atom ein $B$-Atom ist und umgekehrt, kann der Hüpfterm im Hamiltonian mit den Nächsten-Nachbar-Vektoren $\vdelta$ aus Gleichung \eqref{eq:nn-vektoren} einfach umgeschrieben werden, sodass er die Form
\begin{align}
    \oHNN &= - t \sum_{\vec{R}, \vdelta, \sigma} \oad{\vec{R}, \sigma} \ob{\vec{R} + \vdelta, \sigma} + \hc 
\end{align}
annimmt. 
In den nächsten Schritten wird hieraus eine Gleichung zur Bestimmung des Ordnungsparameters $\Delta$ hergeleitet.




% \vspace{5cm}

% Das in dieser Arbeit verwendete Modell basiert auf dem Hubbard-Hamiltonian 
% \begin{align}
%     \oH_{\text{Hubbard}} &= U_C \sum_{\vec{R}} \on{\vec{R}, \up} \on{\vec{R}, \down} - t \sum_{\expval{\vec{R},\vec{R'}}, \sigma} \ofed{\vec{R}, \sigma} \ofe{\vec{R'}, \sigma} \, ,
% \end{align}
% dass den in Gleichung~\eqref{eq:TB} eingeführten Tight-Binding Hamiltonian um eine Coulombabstoßung zweier Elektronen unterschiedlichen Spins auf dem gleichen Gitterplatz erweitert. 
% $\on{\alpha} = \ofed{\alpha} \ofe{\alpha}$ ist hier der fermionische Besetzungszahloperator des Zustands $\alpha$ und $U_C$ die (konstante) Coulombwechselwirkung. \\
% Um dieses Modell nun auf Supraleitung anwenden zu können, wird der Hubbard-Hamiltonian so modifiziert, dass es statt der Coulombabstoßung die in der BCS-Theorie grundlegende effektive Elektron-Elektron-Anziehung beschreibt, sowie auch die Untergitterstruktur des Bienenwabengitters berücksichtigt.
% Der modifizierte Hamiltonian lautet dann
% \begin{align}
%     \oH = - t \sum_{\expval{\vec{R}, \vec{R'}}, \sigma} \sum_{\oC, \oC'}\ocd{\vec{R}, \sigma} \op{C'}^{\pdagg}_{\vec{R'}, \sigma} - U \sum_{\vec{R}} \sum_\oC \on{\oC, \vec{R}, \up} \on{\oC, \vec{R},\down}.
% \end{align}
% wobei nun $U > 0$ die effektive (konstante) Anziehungsstärke beschreibt und daher mit einem Minuszeichen versehen wird, aber für Einelektronenenergien $\epsilon_{\vk}$ außerhalb von 
% \begin{align}
%     E_\text{F} - \hbar \omega_{\text{D}} \leq \epsilon_\vk \leq E_\text{F} + \hbar \omega_\text{D} \label{eq:cutoff}
% \end{align}
% null ist.
% Die allgemeinen fermionischen Operatoren $\ofed{}/\ofe{}$ werden durch Untergitter-spezifische Operatoren $\oad{}/\oa{}$ und $\obd{}/\ob{}$ ersetzt, die nur Elektronen auf dem jeweiligen Untergitter $A$ oder $B$ erzeugen, bzw.~vernichten. 
% % $\oC{}$ ist ein Platzhalter für entweder $\oA{}$ oder $\oB{}$ und 
% Die Summe über $\oC$ bedeutet, dass jeweils über $\oC = \oA$ und $\oC = \oB$ summiert wird. \\

% Da im Bienenwabengitter der nächste Nachbar von jedem $A$-Atom ein $B$-Atom ist und umgekehrt, kann der Hüpfterm im Hamiltonian mit den Nächsten-Nachbar-Vektoren $\vdelta$ aus Gleichung \eqref{eq:nn-vektoren} einfach umgeschrieben werden, sodass der Hamiltonian die Form
% \begin{align}
%     \oH &= \oHNN + \oHWW
% \end{align}
% mit
% \begin{align}
%     \oHNN &= - t \sum_{\vec{R}, \vdelta, \sigma} \oad{\vec{R}, \sigma} \ob{\vec{R} + \vdelta, \sigma} + \hc & & \text{und} & & & \oHWW &= - U \sum_{\vec{R}} \sum_\oC \on{\oC, \vec{R}, \up} \on{\oC, \vec{R},\down}
% \end{align}
% annimmt. 
% In den nächsten Schritten wird hieraus eine Gleichung für den Ordnungsparameter $\Delta$ hergeleitet.




















\section{Fouriertransformation}

% Da die Terme $\oHNN$ und $\oH_\mu$ im Realraum formuliert sind, werden diese in den $\vk$-Raum transformiert um die Translationsinvarianz des Systems auszunutzen, die in der Formulierung von $\oHWW$ bereits eingepflegt ist.
Um die Translationsinvarianz des Systems vollständig auszunutzen, werden die Terme $\oHNN$ und $\oH_\mu$ in den $\vk$-Raum transformiert. 
Dazu werden die fermionischen Operatoren als Fourierreihe der Form
\begin{align}
    \oc{\vec{R}} = \frac{1}{\sqrt{N}} \sum_{\vk} \oc{\vk} e^{\im \vk \cdot \vec{R}} & & \text{und} & & \ocd{\vec{R}} = \frac{1}{\sqrt{N}} \sum_{\vk} \ocd{\vk} e^{-\im \vk \cdot \vec{R}}
\end{align}
geschrieben und in die beiden Hamiltonians respektiv eingesetzt. 
Die Trafo wird mit $\sqrt{N}$ normiert, wobei $N$ für die Anzahl an Gitterplätzen steht.
Der Nächste-Nachbar-Term transformiert damit zu
\begin{align}
    \oHNN &= -t \sum_{\vk,\vq, \sigma} \sum_{\vdelta} \oad{\vq, \sigma}\obk \left(\frac{1}{N}\sum_{\vec{R}} e^{\im \vk \cdot \vdelta} e^{\im \vec{R} \cdot (\vk - \vq)}\right) + \hc\\
    &= \sum_{\vk, \sigma} \underbrace{\left(-t \sum_{\vdelta} e^{\im \vk \cdot \vdelta} \right)}_{:= f(\vk)} \oakd \obk + \hc \\
    &= \sum_{\vk, \sigma} f(\vk) \oakd \obk + \hc \, ,
\end{align}
wobei die Relation $\sum_{\vec{R}} e^{\im \vec{R} \cdot (\vk - \vq)} = N \cdot \delta_{\vk, \vq}$ im ersten Schritt ausgenutzt wird, die gleichzeitig die Impulserhaltung des Systems gewährleistet.
Die Funktion $f(\vk)$ kann für das Bienenwabengitter mit den Nächsten-Nachbar-Vektoren $\vdelta$ aus Gleichung~\eqref{eq:nn-vektoren} explizit als
\begin{align}
    f(\vk) = -t \sum_{\vdelta} e^{\im \vk \cdot \vdelta} = -t \left(\mathrm{exp}(-\im k_x a) + 2 \, \mathrm{exp}\left(\im k_x \frac{a}{2}\right) \cos\left(\frac{k_y a \sqrt{3}}{2}\right)\right) \label{eq:f(k)}
\end{align}
angegeben werden. 

Komplett analog transformiert $\oH_\mu$ zu 
\begin{align}
    \oH_\mu &= -\mu \sum_\oC \sum_{\vk, \vq, \sigma} \ocd{\vk, \sigma} \oc{\vq, \sigma} \underbrace{\left(\frac{1}{N} \sum_{\vec{R}} e^{\im \vec{R} \cdot (\vq - \vk)}\right)}_{=\delta_{\vk, \vq}} \\
    &= - \mu \sum_{\vk, \sigma} \sum_\oC \ocd{\vk, \sigma} \oc{\vk, \sigma} \, ,
\end{align}

womit der Gesamthamiltonian im $\vk$-Raum also 
\begin{align}
    \oH = \sum_{\vk, \sigma} \left(f(\vk) \oakd \obk + \hc - \mu \sum_\oC \ocd{\vk, \sigma} \oc{\vk, \sigma}\right) - \frac{U}{N} \sump{\vk, \vk'} \sum_\oC  \ocdu \ocdd \ocndd \ocndu~\label{eq:H_aMF}
\end{align}
lautet.


% Der Wechselwirkungsteil $\oHWW$ transformiert analog zu
% \begin{align}
%     \oHWW &= -U \sum_\oC \sum_{\vec{m}, \vec{n}, \vec{o}, \vec{p}} \ocd{\vec{m}, \up} \oc{\vec{n}, \up} \ocd{\vec{o}, \down} \oc{\vec{p}, \down} \underbrace{\left(\frac{1}{N^2} \sum_{\vec{R}} e^{\im \vec{R} \cdot \left(\vec{n}+\vec{p} - \vec{m} - \vec{o}\right)}\right)}_{= \frac{1}{N} \delta_{\vec{n} + \vec{p}, \vec{m} +\vec{o}}}\\
%     &= -\frac{U}{N} \sum_\oC \sum_{\vk, \vec{k'}, \vq} \ocd{\vk+\vq, \up} \oc{\vk, \up} \ocd{\vec{k'} - \vq, \down} \oc{\vec{k'}, \down} \\
%     &= - \frac{U}{N} \sum_\oC  \sum_{\vk, \vec{k'}, \vq}\ocdu \ocdd \ocndd \ocndu~\label{eq:Hww} \, ,
% \end{align}
% wobei hier wieder im ersten Schritt Impulserhaltung, also $\vec{n} + \vec{p} =  \vec{m} +\vec{o}$ gewährleistet wird.
% Anhand dieser Bedingung wird die Anzahl an Freiheitsgraden reduziert, indem die neuen Impulse $\vk, \vk', \vq$ eingeführt werden, die zu den alten Impulsen in den Relationen
% \begin{align*}
%     \vec{m} &= \vk + \vq & \vec{o} &= \vk' - \vq \\
%     \vec{n} &= \vk & \vec{p} &= \vk'
% \end{align*}
% stehen. Es ist leicht erkenntlich, dass diese Wahl immernoch $\vec{n} + \vec{p} =  \vec{m} +\vec{o}$ erfüllt.
% Vom zweiten auf den dritten Schritt werden die Antikommutatorrelationen $\anticommute{\oc{\alpha}}{\ocd{\beta}} = \delta_{\alpha, \beta}$ und $\anticommute{\oc{\alpha}}{\oc{\beta}} = 0$ angewandt.
% Der Wechselwirkungsterm beschreibt jetzt im Vergleich mit dem BCS-Hamiltonian~\eqref{eq:BCS} die Bildung von Cooper-Paaren mit Singulett-Kopplung aber dafür endlichem Gesamtimpuls $\vec{K} = \vk + \vk'$ und kann nun mit dem zuvor implizit in der Definition von $U$ formulierte Cutoff~\eqref{eq:cutoff} über die Einteilchenenergie $\epsilon_\vk$ explizit formuliert werden, sodass der Hamiltonian die Form
% \begin{align}
%     \oHWW &= - \frac{U}{N} \sum_{\vk, \vk', \vq} \sum_\oC \ocdu \ocdd \ocndd \ocndu \Theta(\hbar \omega_D - \abs{\epsilon_\vk - E_\text{F}}) \Theta(\hbar \omega_D - \abs{\epsilon_{\vk'} - E_\text{F}}) 
% \end{align}
% bekommt.
% % wobei hier der zuvor implizit formulierte Cutoff über die Einteilchenenergie $\epsilon_\vk$ nun explizit formuliert ist.
% In dem Rest der Arbeit wird die Fermienergie $E_\text{F}$ auf den Energienullpunkt gesetzt und die Schreibweise 
% \begin{align}
%     \sump{\vk} &= \sum_{\vk} \Theta(\hbar \omega_D - \abs{\epsilon_\vk - E_\text{F}}) = \sum_{\vk} \Theta(\hbar \omega_D - \abs{\epsilon_\vk})
% \end{align}
% verwendet, womit der Gesamthamiltonian im $\vk$-Raum 
% \begin{align}
%     \oH &= \sum_{\vk, \sigma} \left(f(\vk) \oakd \obk + \hc \right) - \frac{U}{N} \sump{\vk, \vec{k'}} \sum_{\vq,\oC} \ocdu \ocdd \ocndd \ocndu \label{eq:H_vor_MF}
% \end{align}
% lautet.
















\section{Mean-Field-Näherung}

% Da das Hubbard-Modell selbst im einfachsten 1-Band-Fall nicht für Dimension $1< d < \infty$ exakt gelöst werden kann (vgl.~\cite{hanisch_lattice_1997}), muss an dieser Stelle eine Näherung zur Vereinfachung des Wechselwirkungsteils aus Gleichung~\eqref{eq:H_vor_MF} angewandt werden. 

% In Anhang~\ref{anang_a} wird gezeigt, dass 

% Da das Hubbard-Modell für Dimensionen $1 < d < \infty$ noch nicht exakt gelöst werden kann (vgl.~\cite{hanisch_lattice_1997}), wird an dieser Stelle eine Näherung zur Vereinfachung des Wechselwirkungsteils aus Gleichung~\eqref{eq:H_vor_MF} angewandt.
% Wie in der BCS-Theorie üblich wird hier eine Mean-Field-Näherung durchgeführt, welche für supraleitende Phasenübergänge im Vergleich zu anderen Phasenübergängen gut die experimentell gemessenen Abhängigkeiten des Ordnungsparameters von der Temperatur reproduziert~\cite[234]{czycholl_theoretische_2017}.\\
Um den Wechselwirkungsterm aus~\eqref{eq:H_aMF} mathematisch einfacher behandeln zu können, wird an dieser Stelle eine Mean-Field-Näherung angewandt. 
Diese Näherung ist insofern gerechtfertigt, dass die exakte Lösung im thermodynamischen Limes $N \rightarrow \infty$ genau mit der Mean-Field Lösung zusammenfällt~\cite{bogolyubov_asymptotically_1961,roman_large-n_2002}.
Ebenso kann begründet werden, dass aufgrund der makroskopischen Effekte bei der Bildung des Cooper-Paar-Kondensats sehr viele Teilchen beteiligt sind und dadurch quantenmechanische Fluktuationen in guter Näherung vernachlässigt werden können~\cite[50]{tinkham_introduction_2004}.
Da sich in dieser Arbeit sowieso nur auf den thermodynamischen Limes beschränkt wird, ist die Mean-Field-Näherung hier also eine gute Approximation.



Als Vorbereitung für die Näherung wird der quartische Term aus Gleichung~\eqref{eq:H_aMF} über das Wick-Theorem~\cite{wick_evaluation_1950,kehrein_flow_2006} wie folgt umgeschrieben:
% \ocd{\vk', \up} \ocd{-\vk', \down} \oc{-\vk, \down} \oc{\vk, \up} &= \, : \ocd{\vk', \up} \ocd{-\vk', \down} \oc{-\vk, \down} \oc{\vk, \up}: \label{eq:2TT}\\
% &+ \expval{\ocd{\vk', \up} \ocd{-\vk', \down}} :\oc{-\vk, \down} \oc{\vk, \up}: + \expval{\oc{-\vk, \down} \oc{\vk, \up}} :\ocd{\vk', \up} \ocd{-\vk', \down}: \label{eq:anomal} \\
% &- \expval{\ocd{\vk', \up} \oc{-\vk, \down}} :\ocd{-\vk', \down} \oc{\vk, \up}: - \expval{\ocd{-\vk', \down} \oc{\vk, \up}} :\ocd{\vk', \up} \oc{-\vk, \down}:\label{eq:fock}\\
% &+ \expval{\ocd{\vk', \up} \oc{\vk, \up}} :\ocd{-\vk', \down} \oc{-\vk, \down}: + \expval{\ocd{-\vk', \down} \oc{-\vk, \down}} :\ocd{\vk', \up} \oc{\vk, \up}:\label{eq:hartree}\\
% % &+ \expval{\ocd{\vk', \up} \oc{\vk, \up}} \expval{\ocd{-\vk', \down} \oc{-\vk, \down}} + \expval{\ocd{\vk', \up} \ocd{-\vk', \down}} \expval{\oc{-\vk, \down} \oc{\vk, \up}} \, ,
\begin{align}
    \ocdu \ocdd \ocndd \ocndu &= \, :\ocdu \ocdd \ocndd \ocndu:\label{eq:2TT}\\
    &+ \expval{\ocdu \ocdd} :\ocndd \ocndu: + \expval{\ocndd \ocndu} :\ocdu \ocdd:\label{eq:anomal}\\
    &- \expval{\ocdu \ocndd} :\ocdd \ocndu: - \expval{\ocdd \ocndu} :\ocdu \ocndd:\label{eq:fock}\\
    &+ \expval{\ocdu \ocndu} :\ocdd \ocndd: + \expval{\ocdd \ocndd} :\ocdu \ocndu:\label{eq:hartree}\\
    &+ \text{const.} \, ,
\end{align}
wobei ``$: \dots :$'' für die Normalordnung der Operatoren steht.
Im Folgenden wird diese allerdings vernachlässigt, da für die Normalordnung von bilinearen Termen $: \op{c}\op{c}: = \op{c}\op{c} - \expval{\op{c} \op{c}} $ gilt~\cite{kehrein_flow_2006} und $\expval{\op{c} \op{c}}$ einfach in die Konstante absorbiert werden kann.
Der bis dahin exakte Term wird nun genähert, indem der normalgeordnete Zweiteilchenterm~\eqref{eq:2TT} als klein angenommen und dementsprechend vernachlässigt wird. 
Dies ist die bekannte Mean-Field-Näherung, die ein Vielteilchenproblem auf ein effektives Einteilchenproblem reduziert~\cite[66]{bruus_manybody_2004}.
Die Konstante ist ebenfalls vernachlässigbar, da sie nur eine Verschiebung der Gesamtenergie bewirkt und daher für die Dynamik des Systems irrelevant ist. \\

Von den verbleibenden Einteilchentermen lassen sich die Beiträge~\eqref{eq:hartree} und~\eqref{eq:fock} den klassichen Hartree- und Fock-Termen der Hartree-Fock-Näherung zuordnen~\cite{czycholl_theoretische_2016,bruus_manybody_2004}: 
% Der Hartree-Term~\eqref{eq:hartree} koppelt mit Erwartungswerten $\expval{\ocdu \ocndu}$ bzw. $\expval{\ocdd \ocndd}$ aufgrund der angenommenen Impulserhaltung des Systems (also $\vq = 0$) die Besetzungszahlen der Elektronen mit gleichem Spin.
Der Fock-Term~\eqref{eq:fock} beschreibt hier über die Erwartungswerte $\expval{\ocdu \ocdd}$ und $\expval{\ocdd \ocndu}$ eine Austauschkopplung.
Da ein spinerhaltendes System angenommen wird, sind diese Erwartungswerte allerdings null und der Fock-Term fällt weg.
Der Hartree-Term~\eqref{eq:hartree} beschreibt hier eine direkte Kopplung der Elektronendichte gleicher Spins mit den mittleren Feldern entgegengesetzter Spins $\expval{\ocdu \ocndu}$ bzw. $\expval{\ocdd \ocndd}$. 
Aufgrund der angenommenen Impulserhaltung muss hier $\vk=\vk'$ gelten.
Obwohl er durch diese Einschränkung nicht direkt verschwindet, kann der Hartree-Term aufgrund dessen ebenfalls für die weitere Rechnung vernachlässigt werden.
Dies ist erkenntlich, wenn er in die Wechselwirkung eingesetzt wird und damit den Term
\begin{align}
        -\frac{U}{N} \sum_\oC \sump{\vk} \left[\expval{\ocd{\vk, \up} \oc{\vk, \up}} \ocd{-\vk, \down} \ocndd + \expval{\ocd{-\vk, \down} \ocndd} \ocd{\vk, \up} \ocndu\right]
\end{align}
erzeugt. 
Für $N \rightarrow \infty$ bleibt dieser Term aufgrund der Vorfaktors $\frac{1}{N}$ konstant, während die restlichen Terme im Hamiltonian mit $N$ wachsen.
Der durch den Hartree-Term entstehendende Beitrag geht also im Vergleich gegen null.

Zusätzlich zu den bekannten Hartree- und Fock-Termen gibt es mit~\eqref{eq:anomal} einen weiteren, für die BCS-Theorie zentralen Beitrag: 
Er enthält die anomale Erwartungswerte $\expval{\ocdu \ocdd}$ bzw. $\expval{\ocndd \ocndu}$, also Erwartungswerte der Paarerzeugung(-vernichtung) von Cooper-Paaren mit dem Gesamtimpuls $\vec{K} = 0$. 
Da diese Erwartungswerte die Teilchenzahlerhaltung verletzten, sind sie in der normalleitenden Phase des Systems immer null, für die supraleitende Phase allerdings aufgrund der makroskopischen Kohärenz der Cooper-Paare endlich.
Sie fungieren also als Ordnungsparameter der supraleitenden Phase.
% Da in dieser Arbeit $s$-wave, also konventionelle Supraleitung nach der BCS-Theorie betrachtet wird, muss außerdem nach dem Argument aus Abbildung~\ref{fig:s_wave} $\vk = -\vk$ gelten, damit der Gesamtimpuls der Cooper-Paare verschwindet.
% Diese erhalten zwar die Teilchenzahl nicht und würden im normalleitenden Fall daher wegfallen, sind aber natürlich für die Ausbildung einer supraleitenden Phase von zentraler Bedeutung und verschwinden dort keinesfalls.
% Diese Erwartungswerte fungieren also als Ordnungsparameter der supraleitenden Phase.

% Wird die diskutierte Mean-Field-Entwicklung nun in den Wechselwirkungs-Hamilton-Operator eingesetzt, muss dieser also mit einem effektiven Einteilchen-Hamilton-Operator
% \begin{align}
%     \oHWWeff &= \oH_1 + \oH_2 \, ,
% \end{align}
% bestehend aus den Teilen $\oH_1$ und $\oH_2$ für den anomalen BCS- und Hartree-Term, ersetzt werden.


Wird die diskutierte Mean-Field-Näherung nun in den Wechselwirkungs-Hamilton-Operator eingesetzt, wird dieser zu einem effektiven Einteilchen-Hamiltonian der Form
\begin{align}
    \oHWWeff &= -\frac{U}{N} \sum_\oC \sump{\vk, \vk'}\left[\ocd{\vk, \up} \ocd{-\vk, \down} \expval{\oc{-\vk', \down} \oc{\vk', \up}} + \ocndd \ocndu \expval{\ocdu \ocdd}\right] \, ,
\end{align}
wobei im ersten Teil der Summe $\vk$ und $\vk'$ per Indextrafo im Vergleich zu~\eqref{eq:anomal} vertauscht wurden.
% Im nächsten Schritt wird nun die Indextrafo $\vec{k'} = \vec{k + q}$ und gleichzeitig der Indextausch $\vec{k'} \leftrightarrow \vk $ im ersten Teil der Summe durchgeführt, sodass sich unter Einführung der Ordnungsparameter
Unter Einführung der Ordnungsparameter
\begin{align}
    \Delta_\oC &= \frac{U}{N} \sump{\vec{k'}} \expval{\oc{-\vk, \down} \ocndu} & & \text{und} & \Delta_{\oC}^{*} &= \frac{U}{N} \sump{\vec{k'}} \expval{\ocdup, \ocddown}~\label{eq:def_delta} \,
\end{align}
kann der Hamiltonian in die kompakte Form
\begin{align}
    \oHWWeff &= - \sum_\oC \Delta_\oC \sump{\vk} \ocdup \ocddown + \hc
\end{align}
gebracht werden.
% Analog zum chemischen Potential kann hier ebenfalls aufgrund der Untergittersymmetrie angenommen werden, dass die Ordnungsparameter unabhängig von der Wahl des Untergitters sind und daher die Relation $\Delta_{\oA} = \Delta_{\oB} \equiv \Delta$ gelten muss. 
Aufgrund der Untergittersymmetrie wird für den Ordnungsparameter $\Delta_{\oA} = \Delta_{\oB} \equiv \Delta$ angenommen.
Zusätzlich kann $\Delta$ aufgrund der Eichsymmetrie des Systems der Einfachheit halber reel gewählt werden~\cite[272]{altland_condensed_2012}, womit sich 
$\oHWWeff$ weiter zu 
\begin{align}
    \oHWWeff &= - \Delta \left(\sum_\oC \sump{\vk} \ocdup \ocddown + \hc\right)
\end{align}
vereinfacht. 
Insgesamt ergibt sich damit der effektive Ein-Teilchen-Hamiltonoperator 
\begin{align}
    \oHeff &= \oHNN + \oH_\mu + \oHWWeff \\
    % &= \sum_{\vec{\tilde{k}}}\left[\left(f(\vec{\tilde{k}}) \op{A}_{\vec{\tilde{k}}, \up}^{\dagg} \op{B}_{\vec{\tilde{k}}, \down}^{\pdagg} + \hc \right) - \mu \sum_\oC \ocd{\vec{\tilde{k}}, \sigma} \oc{\vec{\tilde{k}}, \sigma} \right] \\
    &= \sum_{\vk, \sigma} \left[\left(f(\vk) \oakd \obk +\hc\right) - \mu \left( \oakd \oak + \obkd \obk \right) \right]  \\
    &- \Delta \sump{\vk} \left(\oadup \oaddown + \obdup \obddown + \hc \right) \notag \, .
\end{align}\\
% wobei $\vec{\tilde{k}}$ die Impulse beschreibt, die außerhalb des Cutoffs liegen, also $\abs{\epsilon_{\vec{\tilde{k}}}} > \hbar \omega_D$ erfüllen.
% Um nun Aussagen über den Ordnungsparameter $\Delta$ zu erhalten muss die Selbstkonsistenz der Definition von $\Delta$ ausgenutzt werden. 

Um im Folgenden den Ordnungsparameter $\Delta$ bestimmen zu können, müssen die anomalen Erwartungswerte aus der Definition~\eqref{eq:def_delta} berechnet werden.
Dazu ist es notwendig den Hamiltonian zu diagonalisieren, um die Eigenzustände und Eigenenergien des Systems zu erhalten.
Dabei kann die Summe über die Impulse außerhalb des Cutoffs, für die $\abs{\epsilon_\vk} > \hbar \omega_D$ gilt, vernachlässigt werden, da diese bei der Berechnung von $\Delta$ keinen Beitrag leisten.
Um das Problem der Diagonalisierung zu vereinfachen, werden die Nambu-Spinoren
\begin{align}
    \vopsi &= \begin{pmatrix} \oadup \\  \obdup \\ \oadown \\ \obdown \end{pmatrix} & & \text{und} & & \vopsid = \begin{pmatrix} \oaup \, , & \obup\, , & \oaddown\, , & \obddown\end{pmatrix} \, ,
\end{align}
eingeführt um $\oHeff$ in die Matrix-Form 
\begin{align}
    \oHeff &= \sump{\vk} \vopsid \dul{h} \, \vopsi + \text{const.}
\end{align}
mit 
% \dul{h} = \begin{pmatrix} \mu_{\oA, \up} & -f^*(\vk) & \Delta & 0 \\ -f(\vk) & \mu_{\oB, \up} & 0 & \Delta \\ \Delta & 0 & -\mu_{\oA, \down} & f(-\vk) \\ 0 & \Delta & f^*(-\vk) & -\mu_{\oB, \down} \end{pmatrix}
\begin{align}
    \dul{h} = \begin{pmatrix} \mu & -f^*(\vk) & \Delta & 0 \\ -f(\vk) & \mu & 0 & \Delta \\ \Delta & 0 & -\mu & f(-\vk) \\ 0 & \Delta & f^*(-\vk) & -\mu \end{pmatrix}
\end{align}
zu bringen.
Dazu werden die fermionischen Antikommutatorelationen aus Gleichung~\eqref{eq:anticommute} genutzt, sowie teilweise der Indextausch $\vk \leftrightarrow -\vk$ durchgeführt.
Des Weiteren kann nun aufgrund der Symmetrieeigenschaften von $f(\vk)$~\eqref{eq:f(k)} die Relation $f(-\vk) = f^*(\vk)$ genutzt werden, sodass die Matrix $\dul{h}$ zu
\begin{align}
    \dul{h} = \begin{pmatrix} \mu & -f^*(\vk) & \Delta & 0 \\ -f(\vk) & \mu & 0 & \Delta \\ \Delta & 0 & -\mu & f^*(\vk) \\ 0 & \Delta & f(\vk) & -\mu \end{pmatrix}
\end{align}
wird.
Das Problem der Diagonalisierung reduziert sich damit auf die Diagonalisierung der $4 \times 4$-Matrix $\dul{h}$. \\

% In den folgenden Schritten wird $\dul{h}$ für verschiedene Fälle von $\mu$ diagonalisiert, um mithilfe der Selbstkonsistenzbedingung von $\Delta$ eine Gleichung für den Ordnungsparameter herzuleiten.



% kann der effektive Hamiltonian in die Matrix-Form 
% \begin{align}
%     \oHeff &= \sum_\vk \opsid \dul{h} \, \opsi + \text{const.} & & \text{mit} & & \dul{h} = \begin{pmatrix} \mu_{\oA, \up} & -f^*(\vk) & \Delta & 0 \\ -f(\vk) & \mu_{\oB, \up} & 0 & \Delta \\ \Delta & 0 & -\mu_{\oA, \down} & f(-\vk) \\ 0 & \Delta & f^*(-\vk) & -\mu_{\oB, \down} \end{pmatrix} \, ,
% \end{align}
% umgeschrieben werden. \\
% Des weiteren kann aufgrund der Symmetrieeigenschaften des Bienenwabengitters die Relation $f(-\vk) = f^*(\vk)$ genutzt werden, sodass die Matrix $\dul{h}$ zu
% \begin{align}
%     \dul{h} = \begin{pmatrix} \mu_{\oA, \up} & -f^*(\vk) & \Delta & 0 \\ -f(\vk) & \mu_{\oB, \up} & 0 & \Delta \\ \Delta & 0 & -\mu_{\oA, \down} & f^*(\vk) \\ 0 & \Delta & f(\vk) & -\mu_{\oB, \down} \end{pmatrix}
% \end{align}
% wird. $\dul{h}$ wird in der Literatur auch als Bogoliubov-de-Genne- (BdG)-Matrix bezeichnet.~\cite{quelle}\\






% Die Spinabhängigkeit des chemischen Potentials ist nur relevant, wenn externe Magnetfelder betrachtet werden, die die Spinbesetzung der Gitter beeinflussen, was hier nicht der Fall ist und dementsprechend im Folgenden vernachlässigt wird.



% Die eigentliche Mean-Field Näherung ist es hier nun, den quadrilinearen Term der Entwicklung zu vernachlässigen, sodass nur Hartree- und Fock-mäßige Ein-Teilchen-Wechselwirkungen mit einem gemittelten Feld übrig bleiben. 
% Zusätzlich als besondere Eigenschaft der BCS-Theorie bleibt noch ein anomaler Erwartungswert übrig. 
% Die Konstanten werden hierbei ebenfalls misachtet, die diese nur das allgemeine Energieniveau des Systems verschieben und die Dynamiken des Systems nicht beeinflussen. \\

% In diesem Fall fällt außerdem der Fock Term aus der Expansion heraus, da hier antiparallele Spins miteinander koppeln, was die Drehimpulserhaltung verletzt. Es bleiben also nur Terme mit anomalen Erwartungswerten und Hartree Terme übrig, die im folgenden genauer untersucht werden. \\

% \textbf{Anomaler Erwartungswert} \\
% Der anomale Erwartungswert $\expval{\ocndd \ocndu}$ ist in der normaleitenden Phase des Festkörpers immer null, da dieser Paar-erzeugung und -vernichtung voraussetzt und damit die Teilchenzahlerhaltung verletzt. Gemäß der BCS-Theorie, sind diese Prozesse allerdings der Ursprung dafür, dass sich eine supraleitende Phase ausbilden kann, der Erwartungswert darf hier also nicht pauschal auf null gesetzt werden. Aufgrund der Impulserhaltung muss hier also gelten $\vk = - \vec{k'}$. \\
% Damit ergibt sich in der supraleitende Phase der wechselwirkende Teilhamiltonian
% \begin{align}
%     \oH_1 = -\frac{U}{N} \sum_\oC \sum_{\vk, \vq} \left[\ocdu \ocd{-(\vec{k + q}), \down} \expval{\oc{-\vk, \down} \ocndu} + \expval{\ocdu \ocd{-(\vec{k + q}), \down}} \oc{-\vk, \down} \ocndu\right] \, ,
% \end{align}
% der nun mit der Indextrafo $\vec{k'} = \vec{k + q}$ und gleichzeitigem Indextausch $\vec{k'} \leftrightarrow \vk $ im ersten Teil der Summe zu 
% \begin{align}
%     \oH_1 = - \sum_\oC \Delta_\oC \sum_{\vk} \ocdup \ocddown + \hc
% \end{align}
% mit den Ordnungsparametern
% \begin{align}
%     \Delta_\oC &= \frac{U}{N} \sum_{\vec{k'}} \expval{\oc{-\vk, \down} \ocndu} & & \text{und} & \Delta_{\oC}^{*} &= \frac{U}{N} \sum_{\vec{k'}} \expval{\ocdup, \ocddown}\label{eq:def_delta} \,
% \end{align}
% wird. Da das Bienenwabengitter symmetrisch unter Vertauschung von Gitter $A$ und $B$ ist, kann hier angenommen werden, dass für die Ordnungsparameter $\Delta_{\oA} = \Delta_{\oB} \equiv \Delta$ gilt. Zusätzlich kann die komplexe Phase vom Delta aufgrund der $U(1)$-Eichsymmetrie des Systems auf null gesetzt werden~\cite[272]{altland_condensed_2012}, also ist es hier kein Problem $\Delta$ der Einfachheit halber reel zu wählen. Damit kann dann $\oH_1$ weiter zu 
% \begin{align}
%     \oH_1 &= - \Delta \left(\sum_\oC \sum_\vk \ocdup \ocddown + \hc\right)
% \end{align}
% vereinfacht werden. \\

% \textbf{Hartree-Term}\\
% Die Erwartungswerte $\expval{\ocdu \ocndu}$ und $\expval{\ocdu \ocndu}$ der beiden Hartree-Terme verletzen für $\vq \neq 0$ die Impulserhaltung der Systems. Darum muss hier also nur der Beitrag der Hartree-Terme für $\vq = 0$ betrachtet werden, denn für endliche Werte von $\vq$ sind die Erwartungswerte null. \\
% Damit ergibt sich der Hartree-Hamiltonian 
% \begin{align}
%     \oH_2 &= -\frac{U}{N} \sum_\oC \left[ \sum_{\vk, \vec{k'}} \left(\ocdup \ocup \expval{\ocd{\vec{k'}, \down} \oc{\vec{k'}, \down}} + \ocd{\vk, \down} \oc{\vk, \down} \expval{ \ocd{\vec{k', \up}} \oc{\vec{k'}, \up} } \right) \right] \, ,
% \end{align}
% welcher unter der Einführung der spin- und gitter-abhängigen Besetzungszahl 
% \begin{equation}
%     n_{\oC, \sigma} = \frac{1}{N} \sum_{\vk'} \expval{\ockd, \ock}
% \end{equation}
% und dem damit assozierten chemischen Potential 
% \begin{equation}
%     \mu_{\oC, \sigma} = U n_{\oC, \overbar{\sigma}}
% \end{equation}
% schließlich in die Form 
% \begin{align}
%     \oH_2 &= - \sum_\oC \sum_{\vk, \sigma} \mu_{\oC, \sigma} \on{\oC, \vk, \sigma} 
% \end{align}
% gebracht werden kann. \\
% Die bei der Definition des chemischen Potentials verwendete Schreibweise $\overbar{\sigma}$ soll hier zum Ausdruck bringen, dass die Besetzungszahl des gegensätzlichen Spins das chemische Potential erzeugt. Diese Feinheit ist allerdings nur bei Betrachtung spinabhängiger Störungen, wie externen Magnetfeldern relevant und ist für das hier betrachtete Modell unerheblich.
% \\ 

% Insgesamt ergibt sich damit der effektive Ein-Teilchen-Hamiltonoperator 
% \begin{align}
%     \oHeff &= \oHNN + \oH_1 + \oH_2 \\
%     &= \sum_{\vk, \sigma} \left[\left(f(\vk) \oakd \obk +\hc\right) - \mu_{\oA, \sigma} \oakd \oak - \mu_{\oB, \sigma} \obkd \obk\right] \\
%     &- \Delta \sum_{\vk} \left(\oadup \oaddown + \obdup \obddown + \hc \right) \notag \, .
% \end{align}

% In dieser BCS-Näherung ist der Hamiltonian nun exakt lösbar, es verbleibt also nur noch eine Diagonalisierung durchzuführen um Aussagen über die Dynamik des Systems treffen zu können. Der Schlüssel zur Bestimmung des Gap Parameters $\Delta$ und damit der Vorhersage der Existenz einer supraleitenden Phase wird die Ausnutzung der Selbstkonsistenz der Definition von $\Delta$ sein. \\
% Diese Selbstkonsistenzgleichung wird nun in den weiteren Schritten für unterschiedliche chemische Potentiale und damit Füllungen der Gitter hergeleitet.

% \newpage



























\section{Selbstkonsistenz im Spezialfall halber Füllung}

% Um erst einmal ohne größeren Aufwand Aussagen zur Charakteristik der Supraleitung im Bienenwabengitter zu machen, beschäftigen wir uns zuerst mit dem möglichst einfachen Spezialfall, nämlich dem Fall halber Füllung. \\
% In diesem Spezialfall wird die effektive Fermienergie des Systems $E_{F, \eff, \oC} = E_F - \mu_{\oC}$ auf null gesetzt. Wie bereits in \red{ref zu DOS von Bienenwabenitter} gezeigt, sind die Energiezustände um die Energie  $E_\text{symm} = 0$ symmetrisch verteilt, weshalb also bei der Wahl von $E_F = E_\text{symm} = 0$ das Gitter genau zur Hälfte gefüllt ist. \\
Zuerst wird der Spezialfall halber Füllung betrachtet, da in diesem Fall die Selbstkonsistenzgleichung für $\Delta$ am einfachsten hergeleitet werden kann. 
Wie bereits in \autoref{subsec:Bienenwabengitter} argumentiert, liegt halbe Füllung genau dann vor, wenn alle Zustände bis zur Energie $E=0$ gefüllt sind, also die Fermienergie ebenfalls null ist.
% Dies bedeutet konkret, dass das chemische Potential $\mu$ in diesem Spezialfall ebenfalls verschwinden muss, da es die zuvor schon auf null gesetzte Fermienergie $E_F$ verschiebt.
Dies bedeutet, dass das chemische Potential $\mu$ in diesem Spezialfall ebenfalls null sein muss, da es die Verschiebung der Fermienergie vom Nullpunkt beschreibt.

Für $\mu = 0 $ ergibt sich damit die vereinfachte Matrix
\begin{align}
    \dul{h} = \begin{pmatrix} 0 & -f^*(\vk) & \Delta & 0 \\ -f(\vk) & 0 & 0 & \Delta \\ \Delta & 0 & 0 & f^*(\vk) \\ 0 & \Delta & f(\vk) & 0 \end{pmatrix} \, ,
\end{align}
die es nun zu diagonalisieren gilt. 
Gesucht ist also eine unitäre Trafo $\dul{U}$, sodass $\dul{h}_\text{diag} = \dul{U}^\dagger \dul{h} \, \dul{U}$ gilt. 
Diese Trafo wird durch die Lösung des Eigenwertproblems $\dul{h} \, \vec{v}_i = E_i \, \vec{v}_i$ bestimmt, wobei die Eigenvektoren $\vec{v}_i$ die Spalten der Trafo-Matrix $\dul{U}$ bilden.

Die Lösung des Eigenwertproblems liefert 
\begin{align}
    \dul{h}_\text{diag} = \dul{U}^\dagger \dul{h} \, \dul{U} = \diag(E_\vk, E_\vk, - E_\vk, - E_\vk)\, .
\end{align}
mit der zugehörigen unitären Trafo-Matrix\footnote{Es sei hier nur erwähnt, dass bei der Bestimmung der unitären Matrix $\dul{U}$ aufgrund der entarteten Eigenwerte $E_\vk$ darauf geachtet werden muss, dass die berechneten Eigenvektoren orthogonal zueinander sind. Ist dies nicht der Fall muss gegebenfalls ein Gram-Schmidt-Verfahren durchgeführt, um sicherzustellen, dass die Trafo-Matrix auch unitär ist.}
\begin{align}
    \dul{U} = \frac{1}{\sqrt{2}}\begin{pmatrix} 
        1 & 0 & -1 & 0 \\ 
        -\frac{f}{E_\vk} & \frac{\Delta}{E_\vk} & - \frac{f}{E_\vk} & - \frac{\Delta}{E_\vk} \\
        \frac{\Delta}{E_\vk} & \frac{f^{*}}{E_\vk}  & \frac{\Delta}{E_\vk} & - \frac{f^{*}}{E_\vk} \\
        0 & 1 & 0 & 1 \end{pmatrix} \, ,
\end{align}
und den Eigenwerten $E_\vk = \pm \sqrt{\Delta^2 + \abs{f(\vk)}^2}$.
Für $\Delta \rightarrow 0$ werden die Energiebänder des Bienenwabengitters im normalleitenden Grundzustand zurückgewonnen.
Für sie gilt mit der Definition von $f(\vk)$ aus Gleichung~\eqref{eq:f(k)}:
\begin{align}
    \epsilon_{\vk} = \pm \abs{f(\vk)} = \pm t \sqrt{3 + 2 \cos(\sqrt{3} k_y a) + 4 \cos(\sqrt{3} k_y a/2) \cos(3 k_x a/2)}~\label{eq:dispersion} \, ,
\end{align}
und werden in \autoref{subsec:Bienenwabengitter} qualitativ genauer diskutiert.

% Für den Hamiltonian gilt dann 
Der Hamiltonian ergibt sich zu
\begin{align}
    \oHeff &= \sump{\vk} \vopsid \dul{h} \, \vopsi = \sump{\vk} \underbrace{\vopsid \dul{U}}_{=: \vophid} \, \underbrace{\dul{U}^\dagg \dul {h} \, \dul{U}}_{= \dul{h}_\text{diag}} \, \underbrace{\dul{U}^\dagg \vopsi}_{=: \vophi} = \sump{\vk} \vophid \dul{h}_\text{diag} \, \vophi \\
    &= \sump{\vk} E_\vk \left(\ophioned \ophione + \ophitwod \ophitwo - \ophithreed \ophithree - \ophifourd \ophifour \right) \, ,
\end{align}
mit den Quasiteilchen-Spinoren 
\begin{align}
    \vophi = \dul{U}^\dagger \vopsi & & \text{und} & & \vophid = \vopsid \dul{U} \, .
\end{align}
In dieser Quasiteilchenbasis ist der Hamiltonian nun diagonal und formell wechselwirkungsfrei.
Für die Berechnung der Erwartungswerte der Quasiteilchenbesetzungszahlen kann daher die Fermi-Dirac-Funktion $F(E) = (e^{\beta E} + 1)^{-1}$ mit der reziproken Temperatur $\beta = \sfrac{1}{k_B T}$ und Boltzmann-Konstante $k_B$ verwendet werden.
Damit gilt
\begin{align}
    \Phi_{ij} := \expval{\vophi \cdot \vophid}_{ij} = \expval{\ophi_{i, \vk} \ophid_{j, \vk}} &= \expval{1 -\ophid_{j, \vk}  \ophi_{i, \vk}} = (1 - F(E_i))\delta_{ij} \, ,
\end{align}
mit $E_{1,2} = + E_\vk $ und $E_{3,4} = - E_\vk$.
Dies kann nun ausgenutzt werden um die Erwartungswerte der ursprünglichen Fermionenoperatoren $\oc{\vk, \sigma}$ und $\ocd{\vk, \sigma}$ zu berechnen, die in der Definition von $\Delta$~\eqref{eq:def_delta} vorkommen.
Mithilfe der inversen Trafo $\vopsi = \dul{U} \vophi$ ergibt sich
\begin{align}
    \rho_{ij} &:= \expval{\opsi \cdot \opsid}_{ij} = \expval{\dul{U} \ophi \cdot \vophid \dul{U}^\dagger}_{ij} = \left(\dul{U} \expval{\ophi \cdot \vophid} \dul{U}^\dagg\right)_{ij} \, ,
\end{align}
% &= \sum_{kl} {\dul{U}}_{il} \Phi_{lk} {\dul{U}^\dagg}_{kj} 
was unter Ausnutzung von $\Phi_{ij} = (1 - F(E_i))\delta_{ij}$, sowie $(\dul{U})_{ij} = u_{ij}$ und $(\dul{U}^\dagg)_{ij} = u_{ji}^*$ zu  
\begin{align}
    \rho_{ij} = \sum_{lk} u_{il} \, \Phi_{lk} \,  {u^*}_{jk} = \sum_{l} u_{il} \, u_{jl}^* \, (1 - F(E_l))
\end{align}
wird.
Um die Selbstkonsistenzgleichung für $\Delta$ zu erhalten muss aus 
\begin{align}
    \dul{\rho} = \begin{pmatrix}
        \expval{\oadup \oaup} & \expval{\oadup \obup} & \expval{\oadup \oaddown} & \expval{\oadup \obddown} \\
        \expval{\obdup \oaup} & \expval{\obdup \obup} & \expval{\obdup \oaddown} & \expval{\obdup \obddown} \\
        \expval{\oadown \oaup} & \expval{\oadown \obup} & \expval{\oadown \oaddown} & \expval{\oadown \obddown} \\
        \expval{\obdown \oaup} & \expval{\obdown \obup} & \expval{\obdown \oaddown} & \expval{\obdown \obddown} \\
    \end{pmatrix}
\end{align}
einfach der erforderliche Erwartungswert, also entweder $\rho_{13}, \rho_{24}, \rho_{31}$ oder $\rho_{42}$ ausgewählt und in die Definition von $\Delta$ eingesetzt werden. \\
Die Wahl ist allerdings unerheblich für das Ergebnis und mit $\rho_{31}$ 
ergibt sich damit 
\begin{align}
    \Delta &= \frac{U}{N}\sump{\vk} \expval{\oadown \oaup} = \frac{U}{N}\sump{\vk} \rho_{31} = \frac{U}{N}\sump{\vk} \sum_l u_{3l} u_{1l}^* (1 - F(E_l))\label{eq:Delta_rho} \\
    &= \frac{U}{N}\sump{\vk} \frac{\Delta}{2E_\vk}\left(1 - F(E_\vk) - (1 - F(-E_\vk))\right) \\
    &= \frac{U}{2N}\sump{\vk} \frac{\Delta}{E_\vk}\left(F(-E_\vk) - F(E_\vk)\right) \, ,
\end{align}
was mit $F(-x) = 1 - F(x)$ und $\tanh(x) = 1 - 2/(e^{2x} + 1)$  schließlich zu
\begin{align}
    \Delta = \frac{U}{2N} \sump{\vk} \frac{\Delta}{E_\vk} \tanh\left(\frac{\beta E_\vk}{2}\right)~\label{eq:gap_eq_half_filling}
\end{align}
wird.

Dies ist die typische BCS-Selbstkonsistenzgleichung für den Ordnungsparameter $\Delta$~\cite{czycholl_theoretische_2017}. 
Das ist erstaunlich, da von einem bipartiten Gitter mit zwei unterschiedlichen Energiebändern ausgegangen wurde, die klassische BCS-Rechnung jedoch nur ein Energieband berücksichtigt. \\
% Dass hier nur der scheinbare Beitrag des Energiebandes $E_\vk$ auftaucht liegt an der Annahme halber Füllung, bei der sich die Energiebänder $\epsilon_\vk = \pm \abs{f(\vk)}$ nur um ein Vorzeichen unterscheiden und in der Quasiteilchenergie $E_\vk = \sqrt{\Delta^2 + \abs{f(\vk)}^2}$ nur der Betrag relevant ist.
Dies lässt sich als Folge des Spezialfalls halber Füllung verstehen, da die erhaltenen Quasiteilchenenergien $E_\vk = \sqrt{\Delta^2 + \abs{f(\vk)}^2}$ hier unabhängig von dem Vorzeichen von $\epsilon_\pm = \pm \abs{f(\vk)}$, also unabhängig von der Wahl der Energiebänder sind. 
Der Beitrag zur Supraleitung ist damit für beide Energiebänder gleich groß, es muss nicht zwischen Ihnen unterschieden werden.
Effektiv ergibt sich damit also nur die Gleichung für ein Energieband.
% \begin{align}
%     U=\frac{1}{\sqrt2}\begin{pmatrix}
%     -u e^{-i\phi} & v e^{-i\phi} & v e^{-i\phi} & -u e^{-i\phi}\\
%     u & -v & v & -u\\
%     -v e^{-i\phi} & -u e^{-i\phi} & u e^{-i\phi} & v e^{-i\phi}\\
%     v & u & u & v
%     \end{pmatrix}
% \end{align}

Um die Selbstkonsistenzgleichung~\eqref{eq:gap_eq_half_filling} numerisch einfacher behandeln zu können, wird die Summe über $\vk$ mithilfe der Zustandsdichte 
\begin{align}
    \rho(\epsilon) = \frac{1}{N} \sum_\vk \delta(\epsilon - \epsilon_\vk)
\end{align}
in die Integralform
\begin{align}
    \Delta = \frac{U}{2} \int_{-\hbar \omega_D}^{\hbar \omega_D} d\epsilon \, \rho(\epsilon) \frac{\Delta}{\sqrt{\Delta^2 + \epsilon^2}} \tanh\left(\frac{\beta \sqrt{\Delta^2 + \epsilon^2}}{2}\right) \label{eq:gap_eq_half_filling_integral}
\end{align}
umgeschrieben. 
Dabei wird der zuvor in der Schreibweise $\sum'$ implizit dargestellte Energiecutoff $\hbar \omega_D$ nun wieder explizit berücksichtigt.

% Numerische Lösungen bzw. graphische Darstellungen der Selbstkonsistenzgleichung~\eqref{eq:gap_eq_half_filling_integral} werden in Kapitel~\ref{sec:ergebnisse} gezeigt.


\section{Selbstkonsistenz bei beliebiger Füllung}

Nachdem die Selbstkonsistenzgleichung für $\Delta$ im Spezialfall halber Füllung hergeleitet wurde, wird nun der allgemeinere Fall beliebiger Füllung betrachtet.
Es muss diesmal also 
\begin{align}
    \dul{h} = \begin{pmatrix} \mu & -f^*(\vk) & \Delta & 0 \\ -f(\vk) & \mu & 0 & \Delta \\ \Delta & 0 & -\mu & f^*(\vk) \\ 0 & \Delta & f(\vk) & -\mu \end{pmatrix}~\label{eq:h_nothalf}
\end{align}
diagonalisiert werden.
Um dieses Vorhaben zu vereinfachen, kann, wie in Anhang~\ref{sec:anhang_a} gezeigt wird, die Blockstruktur von $\dul{h}$ zur Diagonalisierung ausgenutzt werden.
Es ergibt sich die unitäre Trafo-Matrix 
\begin{align}
    \dul{U} = \frac{1}{2} \begin{pmatrix}
        -\frac{\Delta}{N_1} e^{-\im \varphi} \quad & -\frac{\Delta}{N_2} e^{-\im \varphi} \quad & -\frac{\Delta}{N_3} e^{-\im \varphi} \quad & -\frac{\Delta}{N_4} e^{-\im \varphi} \\[0.3em]
        \frac{\Delta}{N_1} & \frac{\Delta}{N_2} & \frac{\Delta}{N_3} & \frac{\Delta}{N_4} \\[0.3em]
        \left(\frac{-E_+ + \lambda_+}{N_1}\right)e^{-\im \varphi} \quad & \left(\frac{E_+ + \lambda_+}{N_2}\right)e^{-\im \varphi} \quad & \left(\frac{E_- - \lambda_-}{N_3}\right)e^{-\im \varphi} \quad & \left(\frac{-E_- - \lambda_-}{N_4}\right)e^{-\im \varphi} \\[0.3em]
        \frac{E_+ - \lambda_+}{N_1} & \frac{-E_+ - \lambda_+}{N_2} & \frac{E_- - \lambda_-}{N_3} & \frac{-E_- - \lambda_-}{N_4}
    \end{pmatrix}~\label{eq:U_nothalf}
\end{align}
mit 
\begin{align}
    \lambda_\pm &= \mu \pm \abs{f(\vk)}, & E_\pm &= \sqrt{\lambda_{\pm}^2 + \Delta^2}, \\
    N_{1,2} &= \sqrt{E_+(E_+ \mp \lambda_+)}, & N_{3,4} &= \sqrt{E_-(E_- \mp \lambda_-)} \, \\
    \varphi &= \arg(f(\vk)) \, , & e^{-\im \varphi} &= \frac{f^*(\vk)}{\abs{f(\vk)}} \, ,
\end{align}
und der Diagonalmatrix 
\begin{align}
    \dul{h}_\text{diag} = \dul{U}^\dagg \dul{h} \, \dul{U} =\text{diag}(E_+, -E_+, E_-, -E_-)~\label{eq:h_diag_nothalf} \, .
\end{align}
% \begin{pmatrix}
%         E_+ & 0 & 0 & 0 \\ 
%         0 & -E_+ & 0 & 0 \\ 
%         0 & 0 & E_- & 0 \\ 
%         0 & 0 & 0 & -E_-
%     \end{pmatrix}

Damit hat der Hamiltonian in Quasiteilchenbasis die Form
\begin{align}
    \oHeff &= \sump{\vk}E_+\left(\ophioned \ophione - \ophitwod \ophitwo\right) + E_- \left(\ophithreed \ophithree - \ophifourd \ophifour\right)
\end{align}
und mit der zuvor hergeleiteten Rechnungsvorschrift für $\Delta$~\eqref{eq:Delta_rho} folgt 
\begin{align}
    \Delta &= \frac{U}{N} \sump{\vk} \sum_l u_{1l} u_{3l}^* (1 - F(E_l)) \\
    &= \frac{U \Delta}{4N} \sump{\vk}\Bigg[\left(\frac{E_+ - \lambda_+}{{N_1}^2}\right)\left(1 - F(E_+)\right) + \left(\frac{-E_+ - \lambda_+}{{N_2}^2}\right)\left(1 - F(-E_+)\right) \\
    &+ \left(\frac{E_- - \lambda_-}{{N_3}^2}\right)\left(1 - F(E_-)\right) + \left(\frac{-E_- - \lambda_-}{{N_4}^2}\right)\left(1 - F(-E_-)\right)\Bigg] \notag \\
    &= \frac{U \Delta}{4N} \sump{\vk} \left[\frac{F(-E_+) - F(E_+)}{E_+} + \frac{F(-E_-) - F(E_-)}{E_-} \right] \\ \notag \\
    \Longrightarrow  \Delta &= \frac{U}{4N} \sump{\vk} \left[\frac{\Delta}{E_+} \tanh\left(\frac{\beta E_+}{2}\right) + \frac{\Delta}{E_-} \tanh\left(\frac{\beta E_-}{2}\right)\right] \, .
\end{align}

% \sqrt{(\epsilon_{\vk, +} - \mu)^2 + \Delta^2}
% \sqrt{(\epsilon_{\vk, -} - \mu)^2 + \Delta^2}

% Wie zuvor kann die Summe über $\vk$ mithilfe der Zustandsdichte $\rho$ in eine Integralform umgeschrieben werden.
Im Vergleich zum Spezialfall halber Füllung~\eqref{eq:gap_eq_half_filling} haben die beiden Energiebänder $\epsilon_\pm = \pm \abs{f(\vk)}$ nun unterschiedlich große Beiträge zur Supraleitung, da die Quasiteilchenenergien $E_\pm = \sqrt{(\epsilon_{\mp} - \mu)^2 + \Delta^2}$ für $\epsilon_+ \neq \epsilon_-$ nicht mehr die gleichen Beträge liefern.
Bei der Formulierung der Selbstkonsistenz über die Zustandsdichte muss nun also auf formell auf die Beiträge der unterschiedlichen Energiebänder geachtet werden.
Also gilt 
\begin{align}
    \Delta = \frac{U}{4} \int_{\mu -\hbar \omega_D}^{\mu + \hbar \omega_D} \left(\rho_{+}(\epsilon) + \rho_{-}(\epsilon)\right) \frac{\Delta}{\sqrt{(\epsilon - \mu)^2 + \Delta^2}} \tanh\left(\frac{\beta \sqrt{(\epsilon - \mu)^2 + \Delta^2}}{2}\right) \, d\epsilon \, ,
\end{align}
mit den Zustandsdichten $\rho_\pm$, die für die Energiebänder $\epsilon_\pm$ definiert sind.
Da die Fermienergie um $\mu$ verschoben ist, muss hier nun über $\mu \pm \hbar \omega_D$ integriert werden.

Für die Gesamt-Zustandsdichte gilt $\rho(\epsilon) = (\rho_{+}(\epsilon) + \rho_{-}(\epsilon))/2$, sodass die Selbstkonsistenzgleichung für beliebige Füllung in Integralform
\begin{align}
    \Delta = \frac{U}{2} \int_{\mu -\hbar \omega_D}^{\mu + \hbar \omega_D} d\epsilon \, \rho(\epsilon) \frac{\Delta}{\sqrt{(\epsilon - \mu)^2 + \Delta^2}} \tanh\left(\frac{\beta \sqrt{(\epsilon - \mu)^2 + \Delta^2}}{2}\right) \label{eq:gap_eq_arbitrary_filling_integral}
\end{align}
lautet und insbesondere für $\mu = 0$ wieder die Selbstkonsistenzgleichung~\eqref{eq:gap_eq_half_filling_integral} für halbe Füllung zurückliefert. 